\paragraph*{Supervised phase classification.} Carrasquilla and Melko~\cite{carrasquilla2016machine} trained a feed-forward network on Ising configurations labelled by the known $T_c$, evaluated it across the transition, and extracted $T_c$ and critical exponents from the output crossing and a finite-size-scaling collapse; they used convolutional networks for the topological phases of a gauge model. Our variant, in which the training set excludes a window around $T_c$ so that the network never sees the transition region, is a choice of this study. Ch'ng et al.~\cite{chng2016machine} applied convolutional networks to the (quantum) Hubbard model. Suchsland and Wessel~\cite{suchsland2018parameter} study what trained networks encode as a function of their parameters. Mehta et al.~\cite{mehta2018high} provide a tutorial with an Ising notebook that serves as a common benchmark.

\paragraph*{Label-free read-outs.} Wang~\cite{wang2016discovering} used PCA to show that the leading components of Ising configurations carry the order parameter; Hu et al.~\cite{hu2017discovering} examined where such unsupervised read-outs work or break down. van Nieuwenburg et al.~\cite{nieuwenburg2016learning} proposed learning by confusion, in which the transition is the trial boundary at which a classifier trained on deliberately assigned labels is most accurate.

\paragraph*{Sampling and scaling.} Cluster algorithms (Swendsen--Wang~\cite{swendsen1987nonuniversal}, Wolff~\cite{wolff1989collective}) remove most of the critical slowing down of local updates. Finite-size scaling~\cite{fisher1972scaling} and the Binder cumulant~\cite{binder1981finite} give the physical reference read-outs against which a learned $T_c$ should be judged. Other uses of learning in the same area include Boltzmann machines~\cite{torlai2016learning} and learned Monte Carlo proposals~\cite{liu2016self}; they are not part of this study.

\paragraph*{Positioning.} Our contribution is not a new method but a controlled comparison under one metric, with a physical reference, a collapse read-out, per-size biases and sensitivity studies, and with the negative outcomes written down.
